\documentclass[7pt]{article}
%\usetheme{Warsaw}

%\usepackage{beamerinnerthemecircles, beamerouterthemeshadow}
%\usepackage{scrextend}

%\changefontsizes{7pt}


% Standard packages
\usepackage{graphicx}
\usepackage{graphics}

\usepackage{color}
\usepackage[OT4]{fontenc}
\usepackage[utf8]{inputenc}
%\usepackage[latin1]{inputenc}
\usepackage{times}\usepackage{amsmath}
\usepackage[T1]{fontenc}

%\mode<presentation> {
%    \usetheme{Singapore} %Bergen, Berkely, Berlin, Boadilla, CambridgeUS, Darmstadt,
%                          %Frankfurt, Goettingen, Singapore, Warsaw
%    \usecolortheme{default} %beetle, seahorse, wolverine, dolphin, beaver
%}


\title{Wave equation for a 2D membrane (FEM version)}

\definecolor{whi}{rgb}{0.95,0.95,0.95}
% Setup TikZ

%\usepackage{tikz}
%\usetikzlibrary{arrows}
%\tikzstyle{block}=[draw opacity=0.7,line width=1.4cm]


% Author, Title, etc.




% The main document

\begin{document}
\maketitle




\section{the problem}


We solve the wave equation for $\Psi$ its deviation from the equlibrium position at ($x,y$) point,
\begin{equation} \frac{\partial^2 \Psi}{\partial t^2}=v^2\nabla^2\Psi \end{equation}
for a square membrane $(x,y)\in(0,L)\times(0,L)$ using the finite element method. We set the speed of sound equal to $v=1$ and $L=5$.
We will be using the finite element method
 \begin{equation} \Psi(t)=\sum_{i=1}^N c_i(t) h_i(x,y).\end{equation} 
Upon discretization of time we get a formula \begin{equation}\Psi(t+\Delta t)=v^2\Delta t^2 \nabla ^2 \Psi |_t+2\Psi(t)-\Psi(t-\Delta t).\end{equation}
We subsitute the expansion in the shape function basis (1) and project the result on the basis elements. We get a linear system of equations for
the time step for the internal points of the computational box:

\begin{equation}
{\bf O}{\bf c}(t+\Delta t)=\left[-v^2\Delta t^2 {\bf S}+2 {\bf O}\right]{\bf c}(t)-{\bf Oc}(t-\Delta t),
\end{equation}
where ${\bf S}$ is the stiffness matrix (the matrix of the $-\nabla^2$ operator of the 2nd assignement of this course),
and ${\bf O}$ is the overlap matrix.
For the points at the nodes we replace the above equation by $c_k(t+\Delta t)=0$, where $k$ is a node at the boundary.
We add the central point to the boundary and we in each moment in time we set $c_i(t)=\sin\omega t$, where $i$
is the central node of our mesh. 

For the initial condition we take the membrane that is not deformed and not moving $\Psi(0)=\Psi(\Delta t)=0$.




%where ${\bf H}(t)={\bf H}_0+{\bf V} \sin(\omega t)$.

Let us assume that 

\section{Deliverables}
\begin{itemize}
\item find $\Delta t$ for which the scheme is stable.
\item illustrate the membrane vibrations for $\omega=\pi/L$
\item do we reach a stationary state of the vibrations ?
\item what about $\omega=2\pi/L$ and $\omega=\sqrt{2.0}\pi/L$ ?

\end{itemize}
\end{document}




